\documentclass[12pt,oneside,a4paper,article]{abntex2}
\usepackage[utf8]{inputenc} % Codificação do documento
\usepackage[T1]{fontenc}    % Seleção de código de fonte.
\usepackage[brazil]{babel}  % Idioma do documento
\usepackage{graphicx}       % Inclusão de gráficos
\usepackage{tabularx}       % Tabelas avançadas
\usepackage{amsmath}        % Melhorias em matemática
\usepackage{lipsum}         % Geração de texto dummy
\usepackage{authblk}
\usepackage{parskip}

% Configurações específicas do abntex2
% Aqui você pode adicionar configurações específicas, como redefinições de comandos
% ou adições de novos pacotes que são essenciais para o seu documento.

% Carrega o pacote abntex2cite para citações
\usepackage[alf]{abntex2cite} % ou use [num] para citações numéricas

\usepackage[left=3cm,right=2cm,top=3cm,bottom=2cm]{geometry} % Margens
\usepackage{setspace}       % Espaçamento entre linhas
% %\usepackage{natbib}         % Formatação de bibliografia

% Informações de título
\title{\textbf{Documento de Viabilidade de Mangaba's Kitchen}}
\author{Eve
llyn Costa \thanks{evellyn.pedroso@ucsal.edu.br}}
\author{Guilherme Vieira \thanks{guilhermenovais.vieira@ucsal.edu.br}}
\author{Henrique Silva \thanks{henrique.souza@ucsal.edu.br}}
\author[1]{Lucas Santos \thanks{lucassantos.garrido@ucsal.edu.br}}
\author[1]{Maria Luisa \thanks{marialuisacarvalho.silva@ucsal.edu.br}}
\author[1]{Victor Dantas \thanks{victor.dantas@ucsal.edu.br} }
\author[1]{Vinicius Kauan \thanks{viniciuskauan.cruz@ucsal.edu.br} }
\author[1*]{Elton Figueiredo \thanks{elton.figueiredo@pro.ucsal.br}}
\affil{
    Engenharia de Software \par
    Escola de Tecnologias \par
    Universidade Católica do Salvador (UCSAL) \par
    Av. Prof. Pinto de Aguiar, 2589 Pituaçu, CEP: 41740-090 \par
    Salvador/BA, Brasil
}
% \affil[1]{\textit {\{evellyn.pedroso,henrique.souza,lucassantos.garrido,guilhermenovais.vieira\}@ucsal.edu.br}}
% affil[1]{\textit {\{marialuisacarvalho.silva,victor.dantas,viniciuskauan.cruz\}@ucsal.edu.br}}
% \affil[1*]{\textit {\{elton.figueiredo\}@pro.ucsal.edu.br}}




\date{Agosto 2024}

\newpage

\ifthenelse{\equal{\ABNTEXisarticle}{true}}{%
\renewcommand{\maketitlehookb}{}
}{}

% Configurações de aparência do PDF final
% \usepackage{hyperref} % para inserir links
\hypersetup{
    colorlinks=false,       % false: boxed links; true: colored links
    pdfborder={0 0 0},      % remove as bordas ao redor dos links
}

\newpage
\renewcommand*{\Authsep}{, }
\renewcommand*{\Authand}{, }
\renewcommand*{\Authands}{, }
\renewcommand*{\Affilfont}{\normalsize\normalfont}
\renewcommand*{\Authfont}{\bfseries}    % make author names boldface    
\setlength{\affilsep}{2em}   % set the space between author and affiliation

\newsavebox\affbox



\newpage


\begin{document}

\begin{center}
    \includegraphics[width=0.3\textwidth]{imagens-template/ucsal_logo.png} 
\end{center}
{\let\newpage\relax\maketitle}

\newpage
\begin{resumoumacoluna}
\hspace{4.5mm}
O projeto de informatização do Mangaba's Kitchen visa desenvolver um sistema informatizado para o restaurante Mangaba Kitchen, com o objetivo de otimizar o controle de pedidos e estimativas de entrega. A iniciativa busca resolver problemas relacionados à falta de precisão nas previsões de tempo e à dificuldade em gerenciar o fluxo de pedidos, impactando negativamente a experiência dos clientes e a eficiência operacional do restaurante. A solução proposta inclui a criação de um sistema integrado que melhora a comunicação entre garçons e cozinha, oferecendo estimativas de tempo mais precisas e um controle mais eficiente dos pedidos.
O sistema contará com três interfaces principais: uma para o gerente, que gerencia o cadastro de funcionários, pratos e formas de pagamento, além de gerar relatórios de comissão e faturamento; uma para os garçons, que permitirá o gerenciamento de mesas, adição de pedidos e visualização das estimativas de entrega; e uma para os cozinheiros, que possibilitará o gerenciamento dos pedidos em preparação, a visualização das receitas e o controle do tempo de preparo. Cada interface será projetada para atender às necessidades específicas dos usuários e garantir uma operação mais organizada.
Além das funcionalidades principais, como o controle de pedidos e estimativas de tempo, o sistema incluirá recursos de segurança, como logins individuais para os funcionários, e relatórios detalhados para melhorar a eficiência operacional. A informatização do pagamento servirá apenas para controle, sem transações financeiras diretas. Com essas melhorias, o projeto visa não apenas reduzir erros e atrasos, mas também proporcionar uma experiência mais satisfatória para os clientes e uma gestão interna mais eficiente.
\vspace{\onelineskip}
 
\noindent
\textbf{Palavras-chaves}: Scrum, Sprint, Viabilidade.
\end{resumoumacoluna}

\textual

\vspace{12mm}

\chapter{Introdução}
    \vspace{-6mm}
    \hspace{4.5mm}
O Mangaba Kitchen é um restaurante que se destaca por seu cardápio variado, com refeições completas durante o almoço e especialidades em pizza e hambúrgueres à noite. Com o crescimento das operações e a complexidade crescente no gerenciamento de pedidos, o restaurante identificou a necessidade de uma solução tecnológica para melhorar a eficiência e a precisão no atendimento. Atualmente, a coordenação entre os garçons e a cozinha enfrenta desafios que podem levar a atrasos e erros nos pedidos, impactando diretamente a satisfação dos clientes.
\vspace{12mm}

\newpage
\chapter{Objetivo}
    \vspace{-6mm}
    \hspace{4.5mm}
O principal objetivo do sistema para o Mangaba Kitchen é otimizar a operação do restaurante, melhorando o controle sobre as estimativas de entrega dos pratos e o gerenciamento dos pedidos. Com a crescente complexidade na coordenação entre a equipe de garçons e a cozinha, é fundamental estabelecer um sistema eficiente que não só simplifique a comunicação interna, mas também eleve a qualidade do serviço prestado aos clientes. A implementação de um sistema informatizado visa resolver os problemas atuais e promover uma operação mais organizada e ágil.
\par
Para atingir esse objetivo, o sistema será projetado com várias funcionalidades específicas. Primeiramente, será criado um controle de acesso para funcionários, com logins e senhas individuais, assegurando a segurança e a responsabilidade na operação. Além disso, uma interface para a cozinha permitirá a gestão estruturada dos pedidos, incluindo a capacidade de adicionar pratos à “fila de espera” e ajustar o tempo de preparo conforme necessário. Os cozinheiros também terão acesso às receitas detalhadas dos pratos, o que ajudará a garantir a precisão na preparação e a reduzir erros. Paralelamente, os garçons utilizarão uma interface dedicada para registrar pedidos diretamente, o que facilitará a comunicação do salão e da cozinha.
\par
Além dos aspectos operacionais, o sistema também abordará questões relacionadas à transparência e responsabilização. Será implementado um mecanismo para gerar relatórios detalhados sobre a comissão dos garçons e registrar erros na preparação dos pratos, aplicando desconto financeiro em caso de falhas. Também será criado um sistema de crítica para garantir que os garçons verifiquem restrições alimentares dos clientes, com penalizações em caso de negligência. Esses objetivos visam promover um ambiente de trabalho mais eficiente, seguro e orientado à qualidade, resultando em um melhor atendimento ao cliente e uma operação interna mais fluida e bem organizada.
\vspace{12mm}

\newpage
\chapter{Análise de Viabilidade}
    \vspace{-6mm}
    \hspace{4.5mm}
A viabilidade organizacional e operacional do projeto de informatização do Mangaba's Kitchen é altamente favorável, dada a forma como o sistema está alinhado com os objetivos estratégicos do restaurante e sua integração planejada nas operações diárias. O sistema, que inclui funcionalidades como controle de acesso e gerenciamento de pedidos, foi projetado para melhorar a eficiência e a comunicação interna entre gerentes, garçons e cozinheiros, facilitando o fluxo de trabalho e a coordenação entre as diferentes áreas. A escolha de dispositivos como iPads e uma TV touch, combinada com uma interface intuitiva, garante uma transição suave e a aceitação do sistema pela equipe. A abordagem metodológica Scrum permite um desenvolvimento iterativo, ajustando o sistema com base no feedback dos usuários e nas necessidades emergentes, o que assegura que o sistema continue a ser relevante e eficaz. Esse planejamento detalhado e a adaptação contínua asseguram que a implementação seja harmoniosa, sem interromper o fluxo de trabalho existente, e que a equipe possa rapidamente usufruir das melhorias proporcionadas pelo novo sistema.
\par
A viabilidade econômica do projeto de informatização do Mangaba's Kitchen está intrinsecamente ligada à análise das funcionalidades propostas e ao custo associado ao desenvolvimento do sistema. O sistema detalhado e bem delineado pelo contratante permite uma identificação clara dos requisitos, facilitando a realização de uma avaliação precisa e objetiva. A partir dessa análise, é possível planejar com maior acurácia os recursos necessários, assegurando que o desenvolvimento seja concluído nos prazos e custos estimados. Além disso, a clareza dos objetivos e a definição das funcionalidades essenciais contribuem para uma previsão de manutenção simplificada e eficiente, reduzindo possíveis complicações futuras.
\par
Considerando o esforço de desenvolvimento, que abrange a criação de interfaces específicas para gerentes, garçons e cozinheiros, além de funcionalidades como controle de acesso, gerenciamento de pedidos e visualização de receitas, conclui-se que o projeto é viável economicamente. A estrutura proposta não só se alinha com as necessidades operacionais do restaurante, mas também proporciona uma plataforma escalável, capaz de se adaptar a futuras expansões ou melhorias. A simplicidade na manutenção, derivada da organização e clareza do sistema, reforça ainda mais essa viabilidade, garantindo que o Mangaba Kitchen possa operar com eficiência e agilidade, maximizando o retorno sobre o investimento feito no desenvolvimento dessa solução tecnológica.
\par
A viabilidade técnica do projeto é fortemente sustentada pela clareza e simplicidade dos requisitos estabelecidos, reduzindo a complexidade do desenvolvimento e permitindo uma implementação eficiente. A contratada possui ampla experiência nas ferramentas e tecnologias necessárias para a criação de sistemas de gestão, o que garante uma execução precisa e alinhada às expectativas do cliente. Além disso, a escolha de dispositivos como iPads para os garçons e uma TV touch (sugerida pela Contratada) para os cozinheiros evidencia uma estratégia bem planejada que aproveita as capacidades de hardware já estabelecidas, facilitando a integração e operação do sistema no ambiente do restaurante.
\par
A competência da equipe envolvida no projeto, aliada ao conhecimento profundo das ferramentas utilizadas, assegura que todas as funcionalidades requisitadas serão entregues conforme o cronograma e as especificações. A contratada não apenas domina os aspectos técnicos necessários, mas também compreende as nuances do ambiente de um restaurante, permitindo que o sistema desenvolvido seja robusto, intuitivo e de fácil manutenção. Essa combinação de fatores reforça a confiança na execução do projeto, garantindo que ele será uma solução eficaz e duradoura para as necessidades operacionais do Mangaba's Kitchen.
\vspace{12mm}

\newpage
\chapter{Cronograma}
    \vspace{-6mm}
    \hspace{4.5mm}
A fase de especificação de requisitos foi agendada para ser concluída até o fim do quarto sprint, com prazo total de um mês, onde será produzido os seguintes documentos: problema, objetivo, documento de requisitos, documento de viabilidade, identidade visual e figma do projeto. Após essa fase, a equipe iniciará a fase de projeto, que está agendada para ser concluída até o terceiro sprint, dia 6 de outubro, com prazo total de duas semanas.
\par
Os sprints, que têm a duração de uma semana, eles começaram no dia 25 de agosto, com reuniões Scrum nos dias úteis, para revisar o progresso e ajustar as estratégias conforme necessário. Os sprints continuam regularmente, com o quarto sprint previsto para terminar no dia 15 de setembro, marcando o fim da fase de levantamento de análise de requisitos do sistema. O cronograma inclui ainda um sprint adicional no dia 18 de setembro para refinamento da apresentação dos artefatos para o cliente.
\vspace{12mm}

\newpage
\chapter{Conclusão}
    \vspace{-6mm}
    \hspace{4.5mm}
Portanto, a proposta de informatização do Mangaba Kitchen visa transformar como o restaurante gerencia suas operações diárias, abordando diretamente as dificuldades enfrentadas na coordenação entre a equipe e no controle de pedidos. A introdução de um sistema integrado que inclui interfaces específicas para gerentes, garçons e cozinheiros permitirá uma gestão mais eficiente dos pedidos, um controle aprimorado das estimativas de entrega e uma comunicação mais fluida entre as áreas. Com funcionalidades como cadastro e gerenciamento de funcionários, acompanhamento de pedidos e geração de relatórios financeiros, o sistema visa não apenas otimizar a operação interna, mas também melhorar a experiência do cliente, reduzindo atrasos e erros que impactam a satisfação dos consumidores.
\par
A viabilidade econômica e técnica do projeto está bem sustentada pela clareza dos requisitos e pela experiência da equipe de desenvolvimento, que possui conhecimento profundo das ferramentas e tecnologias necessárias. O cronograma estruturado utilizando o framework Scrum, com fases definidas e sprints regulares, assegura que cada etapa do desenvolvimento será cuidadosamente monitorada e ajustada conforme necessário. A combinação de uma estratégia bem planejada, a escolha adequada de dispositivos e a abordagem iterativa garantem que o sistema será implementado eficientemente, alinhando-se às necessidades operacionais do Mangaba Kitchen e promovendo uma operação mais ágil e eficiente.


% Formatação da bibliografia
% bibliographystyle{plain}
% \newpage
% \bibliography{referencias} % Assume que você tem um arquivo referencias.bib

\end{document}
